% =====================================================================
%  pipeline_changes.tex
%  Documentation of modifications to newFrame.py
%  Selective Betting Project -- Group 11
% =====================================================================
\documentclass[11pt,a4paper]{article}

\usepackage[margin=2.5cm]{geometry}
\usepackage{booktabs}
\usepackage{amsmath}
\usepackage[hidelinks]{hyperref}

% Deliberately no listings/xcolor/parskip: this file is written to compile on
% a minimal TeX installation using core packages only.
\setlength{\parindent}{0pt}
\setlength{\parskip}{0.55\baselineskip}

\newcommand{\code}[1]{\texttt{#1}}

\title{Modifications to the Betting Pipeline\\
       \large What changed in \code{newFrame.py}, and why}
\author{Group 11}
\date{\today}

\begin{document}
\maketitle

% ---------------------------------------------------------------------
\section{Purpose of this document}

This document records every change made to \code{newFrame.py}, the reason
for each one, and its measured effect on the data. It is written to be
readable without reference to the code.

Each change is presented in the same four parts: what the code did before,
why that was a problem, what it does now, and what measurably changed as a
result.

A note on the numbers quoted throughout. Figures described as
\emph{measured} were computed directly from
\code{all-euro-data-2025-2026.csv} as part of making these changes, and can
be reproduced. Figures described as \emph{previously established} are taken
from the project's exploratory analysis as recorded in \code{CLAUDE.md} and
were not re-derived here.

% ---------------------------------------------------------------------
\section{Summary}

\begin{table}[h]
\centering
\small
\begin{tabular}{@{}p{6.1cm}p{4.0cm}p{4.4cm}@{}}
\toprule
\textbf{Change} & \textbf{Type} & \textbf{Measured effect} \\
\midrule
Benchmark price separated from execution price
  & Correctness of the research design
  & Usable matches rise from 2{,}931 to 7{,}454 \\[4pt]
Bets settle at the best available price
  & Realism of the backtest
  & Line-shopping gain now reported separately \\[4pt]
Price columns propagated to the backtest
  & Enabling bug fix
  & Without it the change above cannot take effect \\[4pt]
Rest-days calculation corrected
  & Bug fix
  & 71.2\% of rows had the wrong value \\[4pt]
Empty-split guard
  & Fail-loud safeguard
  & Prevents meaningless output being printed \\[4pt]
Six functions added as stubs
  & Scaffolding
  & No behavioural change yet \\
\bottomrule
\end{tabular}
\caption{Overview of modifications.}
\end{table}

% =====================================================================
\part*{Part I --- Changes to existing code}
\addcontentsline{toc}{part}{Part I --- Changes to existing code}

% ---------------------------------------------------------------------
\section{Separating the benchmark price from the execution price}

\subsection{What the code did before}

A single set of columns, \code{PSH}/\code{PSD}/\code{PSA} (the odds offered
by the bookmaker Pinnacle), was used for two completely different jobs at
once:

\begin{enumerate}
  \item \textbf{As the benchmark.} These odds were converted into implied
        probabilities \code{p\_shin\_H/D/A}. The model's ``edge'' was
        defined as how far its own probability sat above this benchmark.
  \item \textbf{As the execution price.} Every simulated bet was priced and
        settled at these same odds.
\end{enumerate}

\subsection{Why that was a problem}

\paragraph{First problem: the test was circular.}
The strategy's claim is that the model finds value the market has missed.
But ``the market'' and ``the price we are paid at'' were the same number.
We were asking whether the model could beat Pinnacle, while paying it
Pinnacle's own prices, and measuring the gap against Pinnacle's own implied
probabilities.

The plain version: this is like marking your own exam paper. Whenever the
model disagreed with Pinnacle, that disagreement was recorded as edge --- but
the payout for being right was set by the very book we claimed to have
outsmarted. A genuine edge and a modelling error look identical under that
arrangement, because there is no independent reference to separate them.

\paragraph{Second problem: most of the data was being thrown away.}
Pinnacle prices are present for only 39.3\% of matches this season
(measured). Because the pipeline required those columns before a match could
be used at all, roughly three matches in five were discarded before any
modelling began --- not because anything was wrong with them, but because one
particular bookmaker's price was missing.

This matters more than it might appear. The project's central difficulty is
statistical power: with too few bets, a real edge and no edge produce
indistinguishable results. Discarding 60\% of the sample makes that problem
substantially worse for no analytical gain.

\paragraph{Third problem: the Pinnacle figures are not trustworthy here.}
It was previously established that in this season's file, the recorded
Pinnacle price exceeds the recorded market maximum in 26.2\% of matches ---
that is, it claims to be better than the best price any bookmaker offered.
No other bookmaker in the file does this even once. For roughly a quarter of
the sample, the backtest was therefore settling bets at a price that was
never actually available.

\subsection{What the code does now}

The two jobs are separated into two different sets of columns, declared
explicitly at the top of the file:

\begin{quote}\small
\begin{verbatim}
BENCHMARK_COLS = ["AvgCH", "AvgCD", "AvgCA"]   # what we measure against
EXECUTION_COLS = ["MaxH", "MaxD", "MaxA"]      # what we are paid at
REFERENCE_COLS = ["AvgH", "AvgD", "AvgA"]      # for comparison only
\end{verbatim}
\end{quote}

The benchmark is now the \emph{closing average} price across all bookmakers.
This is a better reference than any single book: it is the market's
collective final opinion, and it is independent of the specific price we
choose to bet at.

A helper function, \code{pick\_benchmark\_columns()}, selects these columns
with a fallback chain: closing average, then pre-match average, then
Pinnacle. The fallback exists because closing odds only became available in
2019/20; once the full historical dataset is downloaded, earlier seasons will
need the pre-match average instead.

\subsection{Measured effect}

\begin{table}[h]
\centering
\small
\begin{tabular}{@{}lrr@{}}
\toprule
& \textbf{Coverage} & \textbf{Usable matches} \\
\midrule
Old benchmark (\code{PSH}, Pinnacle)      & 39.3\%  & 2{,}931 \\
New benchmark (\code{AvgCH}, closing avg) & 100.0\% & 7{,}454 \\
\bottomrule
\end{tabular}
\caption{Effect of changing the benchmark price source. Measured on
\code{all-euro-data-2025-2026.csv}.}
\end{table}

The usable sample grows by a factor of 2.5. Nothing about the model changed
to achieve this --- the data was always there, and was being discarded by an
unnecessary requirement.

% ---------------------------------------------------------------------
\section{Betting at the best available price}

\subsection{What the code did before}

Bets were priced and settled at \code{PSH}, as described above.

\subsection{Why we changed it}

If we are simulating a strategy someone could actually run, the price in the
simulation should be a price they could actually get. In practice a bettor
holds accounts at many bookmakers and takes whichever offers the best odds
on the selection they want. This is called \emph{line shopping}, and it is
standard practice.

The \code{Max} columns record exactly that: the highest price offered by any
bookmaker in the file. It is also the only execution price with near-complete
coverage (99.8\%, measured).

\subsection{What the code does now}

All three places where a bet is priced --- the main backtest, the threshold
sweep, and the live bet slate --- now use \code{MaxH}/\code{MaxD}/\code{MaxA}.
This consistency matters: if the sweep and the backtest priced bets
differently, the headline figure would not reconcile with the reported
return, and the discrepancy would be very hard to trace.

\subsection{An honest caveat}

Betting at the maximum available price is optimistic. It assumes accounts at
every bookmaker, and it ignores the fact that the best price often comes with
a low maximum stake. It is the standard convention in this literature, and it
is what the project brief commits to, but it should be stated plainly in the
write-up rather than left implicit. Section~\ref{sec:shopping} describes how
we make its contribution visible.

% ---------------------------------------------------------------------
\section{Making the price columns reach the backtest}

\subsection{The problem}

The pipeline runs in two stages that communicate through a file on disk. The
training stage writes its predictions to a CSV; the backtest reads that CSV
back in.

The list of columns written out included only the Pinnacle prices. The
\code{Max} and \code{Avg} columns were therefore not in the file the backtest
read --- so even if the backtest had asked for them, they would not have been
there.

\subsection{Why this is worth documenting separately}

This is the reason the previous change is not simply a matter of swapping
one column name for another. A change to the backtest alone would have failed
silently or crashed, and the cause would have been several hundred lines away
from the symptom.

It is also a useful general lesson for the write-up: wherever a pipeline
passes data through a file, the column list in that file is a hard boundary.
Anything omitted there is invisible downstream, regardless of what the
earlier code computed.

\subsection{What the code does now}

The output column list includes the execution and reference prices. Pinnacle
columns are still written where present, but only for reference --- nothing
bets at them.

% ---------------------------------------------------------------------
\section{Reporting the line-shopping gain separately}
\label{sec:shopping}

\subsection{The risk we are guarding against}

Switching from an average price to the best available price will improve the
reported return. That improvement is real money, but it is \emph{not evidence
that the model works}. Somebody with no model at all, betting at random,
would also do better at the best price than at the average price.

The danger is obvious once stated: a return that improved purely because we
changed which price we quote could easily be presented --- or read --- as the
model having found an edge.

\subsection{What the code does now}

Every bet is settled twice. The selection and the stake are decided once, at
the best available price, and then held fixed; the same bet is then settled a
second time at the average price. Because only the settlement price differs,
the gap between the two returns isolates the price effect exactly.

The backtest prints both, and labels the difference explicitly as mechanical:

\begin{quote}\small
\begin{verbatim}
Return on Investment : +X.XX%   (settled at Max*)
                       +Y.YY%   (same bets at Avg*)
Line-shopping gain   : +Z.ZZ ROI points, mechanical
\end{verbatim}
\end{quote}

Any reader of the output can now see how much of the return came from
shopping the line and how much, if any, is attributable to the model.

% ---------------------------------------------------------------------
\section{Correcting the rest-days calculation}

\subsection{What the code did before}

One of the model's inputs is how many days of rest each team has had since
its previous match. The calculation grouped matches by home team and measured
the gap between consecutive entries:

\begin{quote}\small
\begin{verbatim}
df.groupby('HomeTeam')['Date'].diff()
\end{verbatim}
\end{quote}

\subsection{Why that was wrong}

Grouping by \code{HomeTeam} only sees a team's \emph{home} matches. The gap it
measures is therefore the time since that team last played \emph{at home},
not the time since it last played at all.

In plain terms: a team that played away on Wednesday and at home on Saturday
was recorded as having had a fortnight's rest, because its previous home
match was two weeks earlier. The away fixture in between was invisible to the
calculation.

This is a genuine bug rather than a design choice. The feature was intended
to capture fatigue, and it was systematically reporting tired teams as fresh.

\subsection{What the code does now}

The calculation was moved onto the long-format table that the pipeline
already builds for rolling goal statistics. In that table each team appears
once per fixture regardless of venue, so consecutive entries for a team are
genuinely its consecutive matches. The corrected values are then joined back
onto the main table.

The 14-day cap and the 14-day default for a team's first appearance were kept,
so that the change shows up as a correction rather than as a rescaling of the
feature.

\subsection{Measured effect}

\begin{table}[h]
\centering
\small
\begin{tabular}{@{}lr@{}}
\toprule
Rows whose rest value changed        & 5{,}316 of 7{,}466 (71.2\%) \\
Changed in the direction of \emph{less} rest & 5{,}316 (all of them) \\
Changed in the direction of \emph{more} rest & 0 \\
Mean rest, before                    & 11.96 days \\
Mean rest, after                     & 7.33 days \\
\bottomrule
\end{tabular}
\caption{Effect of the rest-days correction. Measured on
\code{all-euro-data-2025-2026.csv}.}
\end{table}

The direction of the errors is the key diagnostic. A bug of this kind can
only ever \emph{overstate} rest --- it misses intervening fixtures, and it can
never invent one. Finding that all 5{,}316 changes went one way, and none went
the other, is exactly the signature the diagnosis predicts. Had any row moved
the other way, the fix itself would have been suspect.

A concrete example: Getafe played away at Rayo Vallecano on 2 January 2026 and
at home on 9 January. The original code recorded 14 days of rest for the 9
January fixture; the corrected code records 7.

% ---------------------------------------------------------------------
\section{Failing loudly on empty data splits}

\subsection{The problem}

The training function contains hard-coded lists of which seasons to train,
validate and test on. Those lists name seasons from 2015/16 onwards. The
repository currently holds one season only: 2025/26.

Selecting a season that is not present does not raise an error in pandas. It
returns an empty table. The model would then train on nothing, and the script
would continue to print log-loss figures, bet counts and returns --- all of
them meaningless, and none of them visibly wrong.

\subsection{Why this matters more than it looks}

The failure is silent and the output is plausible. Numbers of roughly the
right shape appear in roughly the right places. This is the most dangerous
class of bug in an analysis pipeline, because nothing prompts anyone to
investigate.

\subsection{What the code does now}

The function checks each split and raises an error naming the seasons
requested and the seasons actually present in the data. The script stops
rather than producing output that cannot be trusted.

Given the current single-season dataset, this check fires immediately. That is
the correct behaviour: the honest state of the project is that the fixed-split
model cannot be run until more data is obtained.

% =====================================================================
\part*{Part II --- New functions added as stubs}
\addcontentsline{toc}{part}{Part II --- New functions added as stubs}

\section{What a stub is, and why we wrote them first}

Six functions were added containing a name, its arguments, a written
explanation, the algorithm in commented pseudocode, and an instruction to
stop if called. None of them contains working code, and none is connected to
the running pipeline.

The purpose is to fix the \emph{interfaces} before writing the
\emph{implementations}: what each piece is called, what it receives, and what
it hands back. These decisions are cheap to change while they are comments
and expensive to change once other code depends on them. Writing them down
first also makes the remaining work, and the order it has to happen in,
visible at a glance.

\section{The six functions}

\subsection{\code{load\_and\_prepare} --- read the data correctly}

The supplied data file is malformed in four ways at once: it uses semicolons
rather than commas, uses a comma as the decimal point (\code{1,44} rather
than \code{1.44}), contains 22 separate embedded header rows, and writes dates
in American month/day/year order.

A separate script, \code{load\_export.py}, already handles all four, but was
never connected to the pipeline. This function is that connection.

Its most important content is a set of assertions: exactly 7{,}466 rows, 22
divisions, a specific date range, no unparsed dates, and specific coverage
percentages for the key odds columns. The reason for asserting rather than
checking informally is that all four faults fail \emph{silently}. Reading the
file carelessly does not raise an error; it quietly discards 573 rows and
turns the odds columns into blanks. The visible symptom is a coverage figure
of 6\% where it should be 99.8\% --- something nobody would notice unless they
happened to look.

\subsection{\code{fetch\_seasons} --- obtain the real dataset}

This downloads twenty seasons of historical data, 2005/06 onwards, from the
original source. That start date is not arbitrary: it is when best-price and
average-price columns first appear, and the best price is now our execution
price.

This function does not currently exist anywhere in the repository, despite
both project documents instructing that it be run.

It is the most important item on the list, because it is what makes a
statistically meaningful result possible at all. With a single season, the
confidence interval around any return remains too wide to distinguish a real
edge from zero, no matter how good the model is. Two later pieces of work
depend on this one entirely.

Note that files from the original source are clean --- ordinary commas, ordinary
dates --- so they must \emph{not} be passed through the loader written for the
malformed export, which would misinterpret their dates.

\subsection{\code{bootstrap\_roi} --- attach uncertainty to every return}

This takes a set of bets and returns not just the return, but a 95\%
confidence interval around it and the probability that the true return is
zero or negative.

The reason is the central discipline of the project: a return quoted without
an interval is not a result. The existing headline figure of $+3.48\%$ on 927
bets carries an interval of $[-2.67\%, +9.55\%]$ --- comfortably including
zero. Reported alone, the point estimate reads as a finding; reported with its
interval, it reads as inconclusive, which is what it is.

One technical point is critical enough to state in the function's own
documentation: the resampling must be done over \emph{matches}, not over
individual bets. Two bets on the same match share an outcome and are
therefore correlated. Resampling bets treats them as independent, which makes
the interval narrower than it should be --- producing false confidence in
precisely the direction we must guard against.

This function is a decision point for the project. If the corrected and
uncorrected results cannot be told apart once intervals are attached, the
correct response is to say so and stop, rather than to build further work on
a result that is not there.

\subsection{\code{run\_walk\_forward} --- test across many seasons}

At present the pipeline trains on one fixed block of seasons and tests on
another. This replaces that with a rolling procedure: for each season in turn,
train on everything up to two years before, validate on the year before, and
test on that season, then move forward one year and repeat.

The reasoning is worth stating carefully, because it is counter-intuitive.
Adding more \emph{training} seasons does not make the result more
statistically reliable --- it may make the model slightly better, but it does
not narrow the confidence interval. Only more \emph{test} data does that.
Rolling the procedure across twenty seasons increases the number of
out-of-sample bets from roughly 900 to roughly 4{,}000, which is the
difference between an inconclusive result and a significant one.

The shrinkage weight is re-estimated on each fold rather than once globally.
Fitting it once assumes that the correct amount of shrinkage has been stable
for twenty years. That is unlikely to be true, and testing it is itself one of
the project's stated research questions.

\subsection{\code{train\_two\_stage} --- stop the model copying the market}

Currently the market's own implied probabilities are among the model's inputs.
The model therefore partly reproduces the market price --- and we then shrink
its output back towards that same price. The deviation being corrected is
partly an echo of the thing it is being measured against, so the correction
is applied twice to the same information.

The fix is to split the model. The first stage sees no market information at
all and forms an independent view. The second stage combines that independent
view with the market price. The disagreement between them is then genuine
rather than circular.

The stub flags one implementation trap explicitly. The second stage must be
trained on the first stage's predictions for data it has \emph{not} seen. If
it is given the first stage's predictions on its own training data, those
predictions will look far more accurate than they really are, the second
stage will learn to over-trust them, and performance will collapse on new
data.

\subsection{\code{check\_non\_vacuousness} --- test whether the correction does anything}

This is a check against the possibility that the project's central
contribution is an illusion.

The method shrinks the model's estimates and then applies a selection
threshold. But shrinking every number by a fixed factor and then applying a
threshold may be arithmetically identical to applying a proportionally
different threshold and not shrinking at all. If that is the case, we have
renamed an axis rather than corrected anything, and the headline figure is an
artefact.

The test is direct. For each threshold used with shrinkage, find the threshold
that selects the same number of bets without shrinkage, and compare the two
sets of bets. If they are essentially the same bets, the correction adds
nothing. If shrinkage genuinely \emph{reorders} which matches look attractive,
the sets will differ, and the correction is doing real work.

This check exists because it is the first thing a sceptical examiner will ask
about, and it is far better to have run it ourselves.

% =====================================================================
\part*{Part III --- Assessing the proposals in \code{n.pdf}}
\addcontentsline{toc}{part}{Part III --- Assessing the proposals in n.pdf}

\section{Why these were assessed rather than adopted}

The document \code{n.pdf} contains five proposals for extending the project.
They were assessed individually against two tests: does the project brief
permit it, and does our data support it. Three of the five failed the first
test outright, being named on the brief's list of things not to do. The
remaining material was reduced to what is actually available and useful.

The reasoning is worth stating because the proposals are not neutral. Four of
the five would make the project sound more sophisticated while weakening the
thing it is arguing. The research claim depends on the pipeline being modular
--- a model, then a calibration step, then a shrinkage step, then a selection
threshold --- because the contribution is specifically located at the
selection step. A single end-to-end network has no selection step to correct.

\begin{table}[h]
\centering
\small
\begin{tabular}{@{}p{5.0cm}p{2.4cm}p{7.0cm}@{}}
\toprule
\textbf{Proposal} & \textbf{Decision} & \textbf{Reason} \\
\midrule
1. End-to-end differentiable model & Rejected
  & Removes the modular selection step the claim depends on; named in the
    brief's exclusions \\[4pt]
2. More features & Partly adopted
  & xG, lineups and press metrics are not in our data; shot counts and market
    signals are \\[4pt]
3. Tabular foundation models (TabNet, FT-Transformer, TabPFN) & Rejected
  & Named individually in the brief's exclusions; evidence cited there is that
    they lose to gradient boosting on this problem \\[4pt]
4. Reinforcement-learning staking agent & Partly adopted
  & The agent is excluded by the brief. The underlying question --- does the
    staking rule matter --- is worth answering \\[4pt]
5. Two-stage model & Already planned
  & Independently specified in the brief; implemented to the brief's
    specification, which differs from \code{n.pdf}'s \\
\bottomrule
\end{tabular}
\caption{Assessment of the five proposals.}
\end{table}

\section{Where proposal 5 and the project brief disagree}

Both documents ask for a two-stage model, but they do not describe the same
model. The brief's version was implemented. Three differences matter:

\begin{itemize}
  \item \code{n.pdf} has the first stage predict physical quantities such as
        expected goals and shot counts. The brief has it predict the three
        match outcome probabilities directly.
  \item \code{n.pdf} feeds the second stage raw odds, naming
        \code{PSH}, \code{MaxH} and \code{AvgH}. The brief feeds it the
        de-vigged probabilities. This difference is not cosmetic: passing the
        raw execution price into the model reintroduces exactly the
        circularity removed in Part~I, because the model would then be
        learning from the price it is about to be paid at.
  \item \code{n.pdf} does not mention out-of-fold prediction. This is the
        detail that determines whether the two-stage model works at all, and
        it is described at length in the stub.
\end{itemize}

\section{What was adopted, and the evidence for it}

\subsection{Market dispersion as an uncertainty measure}

This is the substantive gain from \code{n.pdf}, and it strengthens the
central claim rather than decorating it.

The shrinkage step can vary its strength per match, using what the code calls
variance proxies: inputs that answer ``how much should we trust our own
estimate for this particular match''. The proxies currently used are the rest
days of the two teams. Rest days describe the teams; they say nothing about
our confidence in our own estimate.

Bookmaker disagreement does say something about it. Where the books differ
widely on a match, the true probability is less firmly established, our own
deviation from the market is more likely to be noise, and it should be shrunk
harder. Where the books agree closely, a deviation is more meaningful.

That is the winner's-curse mechanism the project exists to study, expressed as
a measurable quantity. The spread between the best and average price,
$\text{Max}/\text{Avg}-1$, was measured on the 2025/26 data: mean $0.0413$,
standard deviation $0.0273$, tenth percentile $0.0190$, ninetieth percentile
$0.0688$. That is a wide enough range to distinguish confident matches from
uncertain ones.

\subsection{Market movement, and why we expect it to fail}

\code{n.pdf} proposes using the movement between opening and closing prices as
a predictive signal. It is available in our data, and it is legitimate to use
--- closing odds are the price at kickoff, so no future information is
involved.

It was tested before implementing it. Defining drift as
$\log(\text{Avg}) - \log(\text{AvgC})$ on 7{,}454 matches:

\begin{table}[h]
\centering
\small
\begin{tabular}{@{}lrr@{}}
\toprule
\textbf{Correlation of drift with residual against\ldots} & \textbf{$r$} & \textbf{$p$} \\
\midrule
the de-vigged \emph{opening} price & $+0.075$ & $<0.001$ \\
the de-vigged \emph{closing} price & $+0.019$ & $0.095$ \\
\bottomrule
\end{tabular}
\caption{Market drift tested against both price references. Measured.}
\end{table}

The first row is trivially true and should not be mistaken for a finding: the
price moved because the opening price was wrong, so of course the movement
predicts the opening price's error. The second row is the one that matters,
because the closing price is our benchmark, and there the relationship is not
significant.

The plain reading: the market has already absorbed whatever caused the
movement. This mirrors the brief's existing finding that rolling form
correlates $+0.0001$ ($p=0.995$) with the market residual.

It is still worth implementing and reporting. A negative result that was
properly tested is evidence of rigour, and ``we examined market movement and
found it priced'' is a sentence the write-up benefits from being able to make.

\subsection{Shot-based features}

\code{n.pdf} asks for expected goals, expected assists, deep completions and
press intensity. None are present in our data source, and obtaining them means
a second dataset with its own join keys and coverage problems.

What is present, at 92.6\% coverage (measured), is shots and shots on target.
Shots on target is the standard inexpensive proxy for expected goals, and it
is preferable to goals for the same reason expected goals is: goals are a
rare, high-variance event, so a team's recent goal count is a noisy estimate
of how well it has been playing, while its recent shot count is a less noisy
one.

The current feature set contains no shot information whatsoever --- every
rolling feature is built from goals. This is therefore the cheapest available
improvement to the model's inputs.

The expectation should again be tempered. Shot-based form is still form, and
form is priced. The honest hypothesis is that this improves the model's
standalone accuracy while adding little to its \emph{deviation} from the
market, which is the quantity the strategy actually trades on.

One coverage trap was found and is documented in the code: the National League
(\code{EC}) has no match statistics at all, at 0\% against 99.7--100\%
elsewhere. Building these features without noticing would silently drop that
division or fill it with nulls.

\subsection{Staking rule, without the learned agent}

\code{n.pdf} proposes a reinforcement-learning agent that chooses stake sizes.
The brief excludes this explicitly, and there is a methodological reason
beyond the exclusion: an agent that learns staking entangles the decision of
\emph{which} bets to place with the decision of \emph{how much} to place, and
the project's claim is specifically about the first of those.

The underlying question is still worth answering. The backtest currently
stakes a fraction of the Kelly criterion, which makes the stake proportional
to the estimated edge. The model's probability estimate therefore controls
both which bets are selected and how large they are, so an overestimated
probability does damage twice: it gets a poor bet selected, and then sizes it
generously.

Since the claim is that \emph{thresholding} selects for estimation error,
measuring it through a staking rule that amplifies the same error confounds
the two effects. Under flat stakes, any decline in return as the threshold
tightens is attributable to selection alone.

There is a second reason. The brief's power calculation --- a per-bet
profit-and-loss standard deviation of $1.18$, and roughly $2{,}900$ bets
needed to resolve a $3.5\%$ edge --- is quoted at flat stakes. The existing
Kelly-staked results are not directly comparable with it.

The recommendation recorded in the code is that flat stakes become the primary
reported result, with Kelly retained as a robustness check. If the corrected
and uncorrected curves separate under Kelly but not under flat stakes, the
effect lies in the sizing rather than the selection, and the claim as written
would not be supported.

% =====================================================================
\section{What has deliberately not been done}

Two items from the project's task list were left alone, because carrying them
out would have changed reported results in ways not visible in the code
comments:

\begin{itemize}
  \item The flag to exclude the 2025/26 season from the test window. This is
        a scoping decision with a measurable cost in bet count, and it should
        be made explicitly rather than folded into an unrelated set of edits.
  \item Attaching confidence intervals to the returns that already print.
        This depends on \code{bootstrap\_roi}, which is not yet implemented.
        Markers have been left at both places in the code where an interval
        is required.
\end{itemize}

The pipeline does not yet run end to end. Two prerequisites remain outside the
code itself: the \code{lightgbm} library is not installed in the working
environment, and the historical dataset has not been downloaded.

\end{document}
